\documentclass[tikz,border=5pt]{standalone}
\usepackage{fontspec,xeCJK,amsmath,amssymb}
\IfFontExistsTF{Times New Roman}{\setmainfont{Times New Roman}}{\setmainfont{TeX Gyre Termes}}
\IfFontExistsTF{Heiti SC}{\setCJKmainfont[AutoFakeBold=1.5]{Heiti SC}}{\setCJKmainfont[AutoFakeBold=1.5]{FandolHei-Regular}}
\usetikzlibrary{arrows.meta,calc}
\definecolor{blue}{HTML}{2455DE}
\definecolor{green}{HTML}{078565}
\definecolor{orange}{HTML}{D77B00}
\definecolor{red}{HTML}{D53636}
\definecolor{purple}{HTML}{7750C5}
\definecolor{ink}{HTML}{28343C}
\begin{document}
\begin{tikzpicture}[x=1cm,y=1cm,font=\fontsize{9.5}{12}\selectfont,
v/.style={-{Latex[length=2.2mm,width=1.4mm]},line width=.85pt},
axis/.style={-{Latex[length=2mm,width=1.2mm]},ink,line width=.6pt},
aux/.style={ink!45,dashed,line width=.5pt},
box/.style={draw=ink!40,rounded corners=2pt,inner sep=5pt,align=center},
lab/.style={fill=white,inner sep=2pt}]
\path[use as bounding box] (-8,-3.45) rectangle(8,3.45);
\node[font=\bfseries\fontsize{11}{14}\selectfont,text width=15.2cm,align=center] at (0,3.05) {电场线方程使用实际线元：角度增量要乘相应长度};
\begin{scope}[shift={(-5.15,0)}]
\node[text=blue,font=\bfseries] at (0,2.35) {直角坐标};
\draw[axis] (-1.55,-1.0)--(1.65,-1.0);
\draw[axis] (-1.55,-1.0)--(-1.55,1.45);
\node[] at (1.6,-1.25) {$x$};
\node[] at (-1.8,1.4) {$y$};
\draw[v,blue] (-1.15,-0.4)--(0.9,0.85);
\draw[aux] (-1.15,-0.4)--(0.9,-0.4);
\draw[aux] (0.9,-0.4)--(0.9,0.85);
\node[] at (-0.1,-0.7) {$dx$};
\node[] at (1.2,0.1) {$dy$};
\node[] at (-0.9,1.1) {$\odot\,dz$};
\node[box,fill=blue!3,text width=4.1cm] at (0,-1.95) {$\displaystyle\frac{dx}{E_x}=\frac{dy}{E_y}=\frac{dz}{E_z}$};
\end{scope}
\begin{scope}[shift={(0,0)}]
\node[text=green,font=\bfseries] at (0,2.35) {圆柱坐标（俯视）};
\draw[aux] (-.5,0) circle(1.05);
\draw[aux] (-0.5,0)--(0.55,0);
\node[] at (-0.1,-0.3) {$\rho$};
\node[] at (0.55,-0.25) {$P$};
\draw[v,green] (0.55,0)--(1.55,0);
\node[] at (1.12,-0.3) {$d\rho$};
\draw[v,orange] (0.55,0)--(0.55,1.15);
\node[] at (0.9,1.25) {$\rho\,d\phi$};
\node[text=purple] at (-0.8,1.35) {$\odot\,dz$};
\node[box,fill=green!3,text width=4.1cm] at (0,-1.95) {$\displaystyle\frac{d\rho}{E_\rho}=\frac{\rho\,d\phi}{E_\phi}=\frac{dz}{E_z}$};
\end{scope}
\begin{scope}[shift={(5.15,0)}]
\node[text=orange,font=\bfseries] at (0,2.35) {球坐标（子午面）};
\draw[axis] (-1.1,-0.8)--(-1.1,1.5);
\node[] at (-1.3,1.55) {$z$};
\draw[axis] (-1.1,-0.8)--(1.55,-0.8);
\draw[aux] (-1.1,-0.8)--(0.15,0.45);
\node[] at (-0.65,-0.05) {$r$};
\draw[ink!60] (-1.1,.0) arc[start angle=90,end angle=45,radius=.8];
\node[] at (-0.78,0.15) {$\theta$};
\draw[v,orange] (0.15,0.45)--(0.93,1.23);
\node[] at (1.13,1.32) {$dr$};
\draw[v,blue] (0.15,0.45)--(0.95,-0.35);
\node[] at (1.25,-0.45) {$r\,d\theta$};
\node[text=green] at (0.4,1.8) {$\otimes\,r\sin\theta\,d\phi$};
\node[box,fill=orange!3,text width=4.1cm] at (0,-1.95) {$\displaystyle\frac{dr}{E_r}=\frac{r\,d\theta}{E_\theta}=\frac{r\sin\theta\,d\phi}{E_\phi}$};
\end{scope}
\node[text width=15cm,align=center] at (0,-3.02) {各小图显示局部正方向；分母为零时回到参数方程判断，不能直接相除。};
\end{tikzpicture}\end{document}